\subsection{Levers by position on the growth--margin plane}\label{sec:levers-position}
% The lever-table macros (sections/tables/levers_position_macros.tex, generated by code/playbook/run_archetypes.py) are loaded in paper1/p1_preamble.tex.

The envelope results of Section~\ref{sec:model-statics} rank the effects of AI at a single point of the frontier. This subsection states the ranking by position and the one test of it that needs no estimate of the decision weight. It is not a decision rule for a particular firm: the companion paper \citep{murthy2026transformation} turns it into a playbook by position and organizational setup (ownership, readiness, regulated processes, labor law, product model, market structure and owner capability), adds research and development and cost of revenue as levers, and removes levers through gates that hold whatever the position.

Let the firm be at the Rule-of-$X$ optimum for its weight $X>k$, with $\bar\sigma=\infty$, spend $\sigma^*$, excess growth $\Delta^*$, growth $g^*=G_i+\Delta^*$ and growth surplus $\Pi$ (\autoref{prop:score}). Compare four proportional levers at the same proportion $x$: automating the share $x$ of the cost of the growth activities, which multiplies demand pull by $(1-x)^{-\alpha}$; raising demand pull by $x$; cutting a non-discretionary line of size $a$ (general and administrative cost, or maintenance R\&D) by $x$; and cutting human service cost $c_h$ by $x$.

\begin{proposition}[Levers by position]\label{prop:levers-position}
\begin{enumerate}
\item[(a)] To first order the four levers raise $R^*_X$ by $x\sigma^*$, $x\sigma^*/\alpha$, $xa$ and $x(1-\epsilon_p)c_h$. The last two also raise the saturated score $R_\infty$; the first two do not, and vanish as $s\to\bar s$. A point of retention $G_i$ raises $R^*_X$ and $R_\infty$ by $X-k$, and capital for labor raises $R^*_X$ if and only if $da>(g^*+\delta_K)\,dk$.
\item[(b)] Automating the growth activities outranks cutting a line of size $a$ if and only if
\[
X>X^\ddagger_0:=k+\frac{a^{1-\alpha}}{\alpha\ell},
\]
and raising demand pull outranks it if and only if $X>k+(\alpha a)^{1-\alpha}/(\alpha\ell)$. If the firm's observed spend $\hat\sigma$ is optimal at the weight $\hat X$, the first condition is $(X-k)/(\hat X-k)>(a/\hat\sigma)^{1-\alpha}$.
\item[(c)] The growth-engine gains are proportional to $h(s)^{1/(1-\alpha)}$, so a line of size $a$ overtakes them once headroom has fallen by the factor $(a/\sigma^*)^{1-\alpha}$.
\item[(d)] Cutting human service cost outranks a line of size $a$ if and only if $(1-\epsilon_p)c_h>a$; under free entry without a scarce input its long-run effect on the score is zero (\autoref{prop:entry}(b)).
\item[(e)] A plan that lowers $m_0$ by $D$ and raises $G_i$ by $\Delta G$ raises $R^*_X$ by exactly $(X-k)\Delta G-D$.
\item[(f)] Ranked at a screening weight $X^{\rm scr}$ instead of $X$, the growth-side levers of (a), (b) and (e) are understated and the others unchanged, so the ranking reverses whenever $X^{\rm scr}<X^\ddagger_0<X$.
\end{enumerate}
\end{proposition}

\noindent The proof is in Section~\ref{app:proofs} of the Supplementary Appendix.

\smallskip
\emph{What the ranking needs from the data.} The data do not identify the level of the weight $\hat X$ that rationalizes a firm's spend (Section~\ref{sec:empirics}), so the useful form of part (b) is its ratio test, which compares $(X-k)/(\hat X-k)$, the owner's net weight relative to the one its spending reveals, with $(a/\hat\sigma)^{1-\alpha}$ and needs only the observed overhead line, the observed growth spend and $\alpha$. At the revealed weight itself the growth engine outranks overhead exactly when the firm spends more on growth than on the overhead line. The switch weight $X^\ddagger_0=k+(\hat X-k)(a/\hat\sigma)^{1-\alpha}$, which also needs demand pull, its comparison with screening weights in part (f) and the retention rate $G_i=g-\hat\sigma/(\alpha(\hat X-k))$ depend on the level of $\hat X$, and we do not report them.

Table~\ref{tab:levers-position} evaluates the proposition at six positions of the 2023--25 core sample: the cell medians of non-acquisition firm-years in a growth corner ($g$ at or above its upper quartile, $m$ below its lower quartile), a balanced upper-quartile scorer, a mature cash generator (growth below the median, margin in the upper quartile), a sub-40 grower, a stalled firm (both below their medians) and the near-frontier elite of Section~\ref{sec:horizons} (score at or above the 95th percentile). Growth spend $\hat\sigma$ is, on the central basis of Section~\ref{sec:horizons-ranges}, selling and marketing above the maintenance level of \lpMaintBase\% of revenue, with all and none of selling and marketing as bounds, and $\alpha=0.3$, the illustrative value of Section~\ref{sec:model} (the central path of Section~\ref{sec:horizons-ranges} uses \CenAlpha); in the table's last column $\delta$ is the replacement rate $\delta_K$. The table shows only quantities that need no weight: the cell medians, the ratio $(a/\hat\sigma)^{1-\alpha}$ and the lever that ranks first at the position's own weight. No weight is estimated.

\IfFileExists{sections/tables/tab_levers_position.tex}{\begin{table}[t]
\centering
\caption{Levers at six positions of the 2023--25 core sample: the ratio test of Proposition~\ref{prop:levers-position}(b), which needs no growth weight}\label{tab:levers-position}
\scriptsize
\setlength{\tabcolsep}{3pt}
\begin{tabular}{@{}lrrrrrrrrllr@{}}
\toprule
 & & & & & & & & & \multicolumn{2}{c}{First at own weight} & \\
\cmidrule(lr){10-11}
Position & $n$ & $g$ & $m$ & S\&M & $\hat\sigma$ & G\&A & $k$ & $(a/\hat\sigma)^{1-\alpha}$ & central & all S\&M & $g+\delta$\\
\midrule
Growth corner & 28 & 29 & \ensuremath{-}23 & 50 & 25 & 24 & 0.09 & 0.97 & growth engine & growth engine & 0.54 \\
Balanced upper quartile & 43 & 21 & 14 & 26 & 1 & 10 & 0.10 & 3.94 & G\&A cut & growth engine & 0.46 \\
Mature cash generator & 52 & 8 & 20 & 19 & 0 & 10 & 0.19 & -- & G\&A cut & growth engine & 0.33 \\
Sub-40 grower & 48 & 17 & \ensuremath{-}5 & 43 & 18 & 19 & 0.12 & 1.05 & G\&A cut & growth engine & 0.42 \\
Stalled & 84 & 2 & \ensuremath{-}12 & 31 & 6 & 20 & 1.21 & 2.23 & G\&A cut & growth engine & 0.27 \\
Near-frontier elite & 13 & 36 & 22 & 20 & 0 & 15 & 0.04 & -- & G\&A cut & growth engine & 0.61 \\
\bottomrule
\end{tabular}

\smallskip
\begin{minipage}{0.98\linewidth}\scriptsize
Notes: cell medians of non-acquisition core firm-years, fiscal 2023--25, in percent of revenue except $n$, $k$ and the ratio; $m$ is the free-cash-flow margin after stock-based compensation; $\alpha=0.3$. $\hat\sigma$ is growth spend on the central maintenance-netted basis: the selling-and-marketing share less a maintenance base of 25.1\% of revenue, the median S\&M share, in the previous year, of firm-years growing at or below the year's median, and zero below it (the rest of S\&M is a cost line); $a$ is the general and administrative share. $k$ is capital spending and capitalized software over revenue divided by $g$. The ratio $(a/\hat\sigma)^{1-\alpha}$ equals $(X^\ddagger_0-k)/(X-k)$ at the weight $X$ at which the position's spending is optimal: a 10\% cut in general and administrative cost outranks a 10\% saving in the cost of growth at every weight whose distance from $k$ is below the ratio times that $X-k$, and at the position's own weight the cut ranks first if and only if the ratio exceeds one; with no growth spend (--) the cut ranks first at every weight. ``Central'' is the lever that ranks first at the position's own weight on that basis; ``all S\&M'' is the same test with all selling and marketing counted as growth spend, a bound under which the ratio lies between 0.51 and 0.81; with no S\&M counted as growth spend, the other bound, the cut ranks first at every position. No weight is estimated (Section~\ref{sec:empirics}). $g+\delta$ is the opex saved per dollar of compute capital needed to raise the score at $\delta=0.25$. Generated by \texttt{code/playbook/run\_archetypes.py}.
\end{minipage}
\end{table}
}{}

The ranking depends on how much selling spend buys growth. On the central basis and at the weight a position's own spending reveals, cutting a share of general and administrative cost outranks automating the same share of the growth engine at \lpNGACutFirstNet{} of the \lpNPositions{} positions of Table~\ref{tab:levers-position} at the central maintenance level; over the 90\% bootstrap interval of that level the growth engine leads at \StLeverMin{} to \StLeverMax{} positions (Section~\ref{supp:stress} of the Supplementary Appendix). The growth engine leads only in the \MakeLowercase{\lpGEFirstNetList}, whose selling and marketing is far above the maintenance level, and \lpNNoGrowthSpend{} positions (\MakeLowercase{\lpNoGrowthSpendList}) spend nothing above it. Where a position has growth spend, the ratio runs from \lpRatioMin{} to \lpRatioMax; by part (b) the growth engine ranks first for an owner whose net weight exceeds that ratio times the revealed one. If all selling and marketing buys growth, the growth engine leads at \lpNGEFirstAllSM{} of the \lpNPositions{} positions, with ratios from \lpRatioAllSMMin{} to \lpRatioAllSMMax; if none does, at \lpNGEFirstNoSM{} positions. This is a ranking against the growth engine, not of every lever: the companion paper, which also values research and development as a cost line, ranks engineering above both at all \lpNPositions{} positions on the central basis \citep{murthy2026transformation}. Whether a firm ranked at a screening weight should cut overhead first depends on the level of the revealed weight, which the data do not identify (Section~\ref{sec:empirics}): over the range of $X-k$ that the proxies imply, from about \XkLo{} to \XkHi, the switch weight $X^\ddagger_0$ can lie on either side of the screening weights, so the table does not answer that question. The gains from the growth engine are transient and those from overhead durable (part (a)), and capital for labor faces a hurdle that rises with growth (part (a)). Overhead here is general and administrative cost; the table does not rank research and development or cost of revenue.
